\section{Discussion}

\subsection{What ARC contributes}

ARC is best understood as an empirical integration framework. Control-system attachment, secure-base representations, interpersonal risk regulation, dyadic coregulation, and computational attachment models are established precedents \parencite{bowlby1982attachment,kobak1993control,murray2006risk,sbarra2008coregulation,petters2019control,gagliardi2022control,stevens2009dynamic}. The contribution proposed here is to place these ideas into one parsimonious sequence with variables that can be repeatedly measured in adult romantic relationships and with explicit conditions under which components should be rejected.

This positioning deliberately lowers the novelty claim while strengthening falsifiability. The central question is not whether feedback regulation is a new metaphor for attachment. It is whether measured relational discrepancy and feedback provide reproducible information about subsequent action and recovery beyond static attachment dimensions and generic temporal baselines. Throughout this account, control and controller are cybernetic terms for feedback regulation toward a reference condition; they do not denote interpersonal dominance, coercion, or a controlling personality style.

\subsection{Where the existing human evidence stops}

Several pieces of the proposed loop already have human support. Actual--ideal closeness discrepancies are associated with relational outcomes \parencite{frost2013closeness,pusch2022closeness,pietras2024closeness}; perceived partner responsiveness fluctuates within people and predicts interpersonal expression and relationship dynamics \parencite{ruan2020responsiveness,segal2016responsiveness}; partner-buffering research further shows that relationship partners can attenuate insecurity-related difficulties \parencite{simpson2014buffering,overall2022buffering}; and attachment orientations moderate within-person affective or dyadic responses to relational cues \parencite{sadikaj2011dynamics,beck2024daily,xie2025interdependence,carnelley2025phubbing}.

What is not yet established is the specific ARC chain linking a momentary, explicitly measured desired--perceived discrepancy to subsequent regulatory action and then to partner feedback and state updating. This distinction prevents adjacent evidence from being treated as direct confirmation.

It also motivates a stricter treatment of the discrepancy construct. Difference scores are compact, but they encode equality and sign constraints on the effects of their components. ARC therefore treats the difference score as a parsimonious hypothesis that must compete against the raw components and asymmetric discrepancy models. A failure of the simple difference representation would revise the mathematics without necessarily eliminating the broader feedback-regulation account.

\subsection{Attachment dimensions are policy moderators, not controller labels}

ARC does not interpret anxiety and avoidance as fixed controller types. The current formulation assigns them a narrower empirical role: they may explain why the same momentary discrepancy is translated into different regulatory channels. This distinction matters because overt approach, interpersonal distancing, and internal disengagement are not interchangeable manifestations of one scalar action tendency. Existing within-person and dyadic evidence supports attachment-related heterogeneity in responses to relational cues, and interactionist theory emphasizes that these effects depend on context and activation conditions; the specific ARC cross-level interactions nevertheless remain untested \parencite{campbell2011interactionist,sadikaj2011dynamics,xie2025interdependence,carnelley2025phubbing}.

This formulation also provides a clear failure mode. If \Anx\ and \Avo\ predict only stable mean differences in behavior and do not alter within-person discrepancy-to-action mappings, ARC should retain them as descriptive attachment dimensions while dropping their proposed dynamic policy role.

\subsection{Positive security is a stronger claim than low insecurity}

The secure-base literature supports positive representational and behavioral aspects of security \parencite{waters2006workingmodels,waters2019securebase,crowell2002securebase}. Separate literatures also link attachment and dyadic coping to recovery from conflict and physiological stress \parencite{powers2006cortisol,meuwly2012cortisol,salvatore2011recovery}. These literatures justify keeping \Sec\ as a prespecified candidate rather than dropping positive security from the research program. What remains uncertain is the bridge between these lines: whether a positive secure-base representation predicts recovery after anxiety and avoidance are already modeled.

H3 therefore intentionally adopts an incremental-validity criterion. A zero-order association between \Sec\ and recovery is not sufficient. Nor is evidence that secure partners recover better when security is operationalized only as low insecurity. The positive-security construct earns promotion to the ARC core only if a separately measured \Sec\ construct improves a preregistered primary recovery outcome beyond \Anx/\Avo, adds out-of-sample predictive value, and the increment is independently replicated or supported by a second preregistered design. This predictive criterion does not by itself establish a causal effect of \Sec. Stronger mechanistic interpretation would require experimental or otherwise identification-oriented evidence, and the recovery outcome must distinguish genuine state/function improvement from mere retreat of the desired relational target.

\subsection{Adult romantic relationships as the empirical target}

Much computational attachment work has been developed around infant--caregiver systems or simulation. Adult romantic relationships add different targets, reciprocal agency, relationship-specific histories, explicit communication, autonomy goals, and partner-specific regulation. Autonomy need satisfaction can coexist with relatedness in close relationships rather than representing its simple opposite \parencite{kluwer2020autonomy}. ARC therefore treats dyadic reciprocity as a substantive empirical issue rather than an automatic extension from infant models.

\subsection{Why the framework is deliberately small}

ARC retains only \Anx\ and \Avo\ as required person-level variables, together with the minimal momentary loop needed to test feedback regulation. \Sec\ is a prespecified candidate for promotion after H3, not a required construct. \MentAff\ and \Flex\ are optional moderators. Relational-control propensity is evaluated in a separate psychometric program, while the recoverability of P/PI/PID and other candidate mechanisms is a separate identifiability question.

These neighboring research programs can inform later extensions of ARC but are not prerequisites for the framework itself. This separation prevents a failure of one speculative component or one mechanistic model from being hidden by reinterpretation of another.

\subsection{Security and disorganization}

Positive secure-base representation and recovery provide a way to test whether security contributes information that cannot be reduced to low insecurity. This is an integration of established secure-base ideas with repeated temporal measurement, not a claim that positive security is newly discovered or already required by ARC.

The disorganization extension is more tentative. Poor recovery, context-residual inconsistency, persistent approach--avoidance conflict, and maladaptive switching are candidate trajectory properties, not definitions of adult disorganization. Their value depends on incremental prediction of independent criteria.

\subsection{Limitations}

The framework has no human validation data. Momentary targets may be reactive to measurement, relationship goals may change rapidly, partner effects create reciprocal causation, and the sampling interval can change the apparent dynamics. Repeated prompts that ask participants to compare perceived and desired relational states may increase discrepancy salience, shift the desired target, or alter regulatory behavior. The pilot therefore treats study-day trends in \(\hat{\mathbf r}_t\), \(\mathbf r^\ast_t\), and \(\mathbf e_t\) as reactivity diagnostics and examines whether primary conclusions are sensitive to temporal-trend adjustment. Such trends cannot uniquely identify measurement reactivity, because genuine relational change is observationally similar, so they are interpreted as a limitation and sensitivity check rather than a correction. Measurement error in \(\hat{\mathbf r}_t\) and \(\mathbf r^\ast_t\) directly contaminates discrepancy. Person-specific processes may also differ enough that group-level coefficients conceal important heterogeneity.

A further limitation is observational equivalence. Multiple computational mechanisms can produce similar action trajectories. ARC therefore prioritizes predictive and discriminating evidence over mechanistic labels.

\subsection{Immediate empirical priority}

The next scientific step is a feasibility pilot centered on H1. If perceived and desired relational states cannot be measured reliably on matched coordinates, or their discrepancy does not predict subsequent action beyond context and action history, the model should be simplified before any large-scale test of moderators or disorganization.
